\section{BA: an inference-time engine for Egyptian Arabic}\label{sec:engine}

\vox already clears every G3 bar for Egyptian Arabic, English and French, so the planned fine-tune was not needed to
list Egyptian, and no Egyptian training data with a licence and provenance we could accept exists beyond two
synthetic voices (NileTTS). \BA therefore keeps the \vox weights unchanged and adds an engine around them. Each
component targets a problem seen in the bake-off; each default was chosen on the tuning sentences only. The
ablations in \cref{sec:results} show which of them move the automatic metrics.

\paragraph{Egyptian number and symbol normaliser}
Digits were the largest error source for every system: on sentences with numbers, CER rose for all of them, and
the listening check noted that ``in most systems the number is not said the Egyptian way''. The normaliser writes
numbers as Cairenes say them, never in MSA: cardinals up to the billions, years, dates, times (\emph{7:30}
$\rightarrow$ \ar{sab3a-wnoss}), prices with currencies, percentages, decimals and simple fractions, phone numbers
digit by digit, ordinals, counted plurals (\emph{4 days} $\rightarrow$ \ar{arba3-ayyam}, with the Egyptian
counted form) and common symbols and abbreviations. Forms follow colloquial usage (\ar{khamastashar} ``fifteen'',
\ar{tultumiyya} ``three hundred''), not MSA (\ar{khamsata3ashar}, \ar{thalathumi2a}). English and French have
their own normalisers with the same coverage.

\paragraph{Loanword policy for code-switched text}
Egyptians mix English words into Arabic sentences. Left in Latin script, \vox reads them with English phonology;
written in Arabic script as Egyptians say them (\emph{meeting} $\rightarrow$ \ar{meeting}), they are read as part of
the Arabic sentence. \BA writes a list of common loanwords in Arabic script, spells short acronyms letter by letter
and leaves other Latin words as they are. On the tuning set this kept intelligibility level (canonical-script CER
\R{cs:keep:cer} kept vs \R{cs:arabic:cer} transliterated) and raised the Egyptian share from \R{cs:keep:c1pct} to
\R{cs:arabic:c1pct}.

\paragraph{Reference transcript}
\vox can condition on the reference audio alone or on the audio and its transcript. In the bake-off the transcript
raised speaker similarity from \R{bk:voxcpm2:ecapa} to \R{bk:voxcpm2-prompt:ecapa} and the Egyptian share from
\R{bk:voxcpm2:c1pct} to \R{bk:voxcpm2-prompt:c1pct}. Users rarely type what they said, so \BA transcribes the
reference with Whisper large-v3-turbo \citep{openai2024whisperturbo}, which on real colloquial Egyptian speech is
close to large-v3 (\R{asr:turbo:cer} vs \R{asr:large:cer} CER) at a fraction of the cost.

\paragraph{Reference preparation}
Uploads are converted to mono 16\,kHz, screened, and cut to the cleanest 6--15\,s window: long pauses are shortened
and the level is normalised. A recording is rejected, with a message the user can act on, if it is too short, too
noisy or clipped. The noise and clipping thresholds were set midway between the worst clean and the best degraded
references in a controlled test, and reject none of the \R{ref:clean:n} real clean clips we checked.

\paragraph{Long text}
\vox degrades on long inputs. \BA splits text into sentences, packs them into chunks of at most
\R{chunk:maxchars} characters without breaking a sentence, generates each chunk, checks its speaking rate and level
against the bands measured on real speech, retries a failing chunk up to twice with a new seed, and joins the chunks
with pauses that depend on the punctuation.

\paragraph{Generation settings and watermark}
Classifier-free guidance 2.0 and 10 diffusion steps, chosen by a six-setting sweep on the tuning sentences. Every
output is watermarked with AudioSeal \citep{sanroman2024audioseal} with a fixed 16-bit message: the mark is computed on
a 16\,kHz copy and added to the 48\,kHz output, leaving content above 8\,kHz untouched (\cref{sec:release}).

\paragraph{What \BA does not change}
No weights are modified. Any gain over \vox therefore comes from what the model is given (text, reference, transcript)
and how its output is checked, which the ablations in \cref{sec:results} separate.
